Optimization pipelines in LLVM are sequences of passes whose effects depend on order and on each other. We ran \NumSweeps{} pipeline variants over \NumKernels{} kernels from four benchmark suites and measured the interactions between pairs of \NumThreePasses{} passes. The main measure is the exact instruction count of the optimized code. Of the \InstInteractionCells{} ordered pairs, \InstInteractionConsistentCells{} have a mean interaction that exceeds two standard errors across kernels. Where a second application of a pass leaves the code unchanged, the self-interaction equals minus the pass's single effect, and this identity holds in all \InstSelfIdentityCases{} such cases. Wall-clock timing resolves few pass interactions: its median noise is \NoiseMedianAbsPct{}\% of the canonical cost, and the \InteractionSignificantCells{} of \InteractionCells{} timing interactions that pass a two-standard-error test lie below the noise floor. Redundant pass executions are common, with a mean ratio of \RpbOverallMean{}. A block-local checker encodes \BlocksEncodablePct{}\% of matched block pairs. Within the encoded names, \NamesRefine{} are proved to refine, and \NamesViolated{} are violated and are candidates for inspection. Of the Pareto-front members, \ParetoMemberUnsatPct{}\% are proved to refine, \ParetoMemberSatPct{}\% have candidate counterexamples, and the rest are not comparable, so the study does not establish that Pareto-optimal schedules refine their canonical builds. The results are bounded by one compiler release, one timing lowering level, and one input per kernel. Of the \OutputsChecked{} sweeps, \OutputsAgreeing{} reproduce the canonical output under the -O0 lowering used for timing. For the \OutputsDiffering{} sweeps that differ, the optimized IR reproduces the canonical output under interpretation and -O2 lowering.
